\section{Limitations}

\paragraph{Character-unit reconstruction.}
The OpenBookQA character floor was added after the original audit identified an evidence hole. It is now recomputed from raw parquet and passes a 12/12 character/token self-test, so we drop the earlier ``unrecomputable'' status. The remaining limitation is narrower: the character values are reconstructed from raw option text rather than retained as per-item fields in the scored records, so item matching to the token leg follows the shared loader and self-test, not a stored character-count column. The convention of excluding the continuation's leading space was explicit; including it yielded the identical 0.363500 value.

\paragraph{One family at one depth cannot identify healing.}
The prospective healed contrast has one non-OLMo family at one absolute depth, leaving Llama-2, Llama-3, and the \texttt{k10}/\texttt{k12}/\texttt{k14} rungs confounded. The observed ladder is a cliff, not a gradient, and we fit no depth curve.

\paragraph{The healing corpora are unmatched and cannot be matched from available data.}
The Qwen3 arm heals on 5.72 epochs of 5.541B SlimPajama tokens, OLMo-2 on 1.0 epoch of 31.7B Dolmino tokens; Qwen3 cannot consume OLMo-2-token Dolmino and raw Dolmino text is on neither disk. Relative depth is untested too: Qwen3's 36 layers to OLMo-2's 32 make \texttt{keep8} retain 22.2\% versus 25.0\%.

\paragraph{The planned causal contrast has dissolved.}
\texttt{H\_heal} is not merely unmeasured: its antecedent required the unhealed Qwen3 \texttt{k8} comparator to sit below an arm-conditional null, and it does not ($-0.139$ points, $p=0.0964$). No amount of further training repairs a comparator failing the criterion. \texttt{H\_family} is equally unavailable because 0/27 cells are significantly below the permutation null. A future read-out can ask only whether the healed arm has material item-level signal.

\paragraph{Regime, pooling, and literature scope.}
Tables mixing truncate-only non-OLMo arms with prune-then-heal OLMo-2 arms cannot identify a family effect. Pooled results average per-item deviations over disjoint benchmarks and must not be read as per-benchmark verdicts. Finally, the Cho et al. overlap audit used arXiv v4 (January 12, 2026); the January 26, 2026 camera-ready could not be diffed, the OpenReview PDF endpoint being inaccessible from the audit network.
